\addcontentsline{toc}{subsubsection}{Solution 54. The Harmonic Involution}

\begin{dlsummary}[title={\bfseries Solution 54. The Harmonic Involution},breakable]

\begin{proof}
The complete family of solutions is
\[
\boxed{f(x)=\frac{c}{x}\qquad(x>0),\quad c>0.}
\]

\paragraph{1. Recover strict monotonicity.}

Rewrite the inequality as
\[
f(x+y)\le\frac{f(x)f(y)}{f(x)+f(y)}<f(x).
\]
Whenever \(b>a>0\), take \(x=a\), \(y=b-a\). Thus
\[
f(b)<f(a),
\]
so \(f\) is strictly decreasing.

\paragraph{2. Obtain an integer scaling bound.}

Taking reciprocals gives
\[
\frac1{f(x+y)}\ge\frac1{f(x)}+\frac1{f(y)}.
\]
Repeated addition therefore yields
\[
\boxed{f(nx)\le\frac{f(x)}n} \tag{1}
\]
for every positive integer \(n\).

\paragraph{3. Make the scaling bound exact.}

Fix \(x>0\) and a positive integer \(n\). From (1),
\[
nf(nx)\le f(x).
\]
Since \(f\) is decreasing,
\[
f\bigl(nf(nx)\bigr)\ge f(f(x))=x. \tag{2}
\]

But applying (1) with input \(f(nx)\) gives
\[
f\bigl(nf(nx)\bigr)
\le\frac{f(f(nx))}{n}
=\frac{nx}{n}=x. \tag{3}
\]
Hence equality holds in (2) and (3). Strict monotonicity makes \(f\) injective, so
\[
nf(nx)=f(x).
\]
Therefore
\[
\boxed{f(nx)=\frac{f(x)}n.} \tag{4}
\]

\paragraph{4. Extend from rationals to all positive reals.}

Put \(c=f(1)>0\). Equation (4) gives
\[
f(1/n)=nc,
\qquad
f(m/n)=\frac{nc}{m}=\frac{c}{m/n}.
\]
Thus \(f(q)=c/q\) for every positive rational \(q\).

For any real \(x>0\), choose positive rationals \(r_k<x<s_k\) with both sequences tending to \(x\). Monotonicity gives
\[
\frac{c}{s_k}=f(s_k)\le f(x)\le f(r_k)=\frac{c}{r_k}.
\]
The bounds converge to \(c/x\), proving \(f(x)=c/x\).

\paragraph{5. Verify every member of the family.}

For any \(c>0\),
\[
f(f(x))=\frac{c}{c/x}=x,
\]
and
\[
f(x+y)\bigl(f(x)+f(y)\bigr)
=\frac{c}{x+y}\left(\frac cx+\frac cy\right)
=\frac{c^2}{xy}
=f(x)f(y).
\]
Thus every \(c>0\) works.
\end{proof}
\end{dlsummary}

\begin{takeaways}
\begin{itemize}
\item \emph{No takeaways were provided in the source draft.}
\end{itemize}
\end{takeaways}
